\documentclass[10pt]{article}
\usepackage[T1]{fontenc}
\usepackage[utf8]{inputenc}
\usepackage[french]{babel}
\usepackage{amsmath,amssymb,mathrsfs,mathpazo}
\usepackage{geometry}
\geometry{a4paper,margin=2cm}
\usepackage{graphicx,array,tabularx,tabularray,colortbl,multirow,multicol,booktabs,enumitem,pifont,ulem,fancybox,diagbox,cancel}
\usepackage{pstricks,pstricks-add,pst-plot,pst-tree,pst-node,pst-text,pst-func,pst-math,pst-bspline,pst-3dplot,pst-eucl,pst-solides3d}
\let\pscalculate\undefined
\usepackage{pst-bezier}
\usepackage{tikz,pgfplots,tkz-tab}
\usetikzlibrary{arrows,arrows.meta,patterns,calc,positioning,trees,decorations.markings,shapes.geometric,shapes.misc,shapes.arrows,intersections,decorations.pathmorphing,decorations.pathreplacing,backgrounds,shadows,babel}
\pgfplotsset{compat=1.18}
\usepgfplotslibrary{fillbetween}
\usepackage[np]{numprint}
\usepackage{listings,makecell,pas-tableur,twoopt}
\let\barre\undefined
\usepackage[pstricks]{tablvar}
\definecolor{darkgreen}{RGB}{0,100,0}
\providecommand{\red}{\color{red}}
\providecommand{\blue}{\color{blue}}
\providecommand{\green}{\color{green}}
\providecommand{\black}{\color{black}}
\providecommand{\white}{\color{white}}
\providecommand{\gray}{\color{gray}}
\providecommand{\cyan}{\color{cyan}}
\providecommand{\magenta}{\color{magenta}}
\providecommand{\yellow}{\color{yellow}}
\providecommand{\violet}{\color{violet}}
\providecommand{\orange}{\color{orange}}
\providecommand{\euro}{\texteuro}
\providecommand{\EUR}{\texteuro}
\providecommand{\ieme}{\textsuperscript{e}}
\providecommand{\degres}{\ensuremath{^\circ}}
\providecommand{\degree}{\ensuremath{^\circ}}
\providecommand{\vv}[1]{\overrightarrow{#1}}
\providecommand{\Coord}[1]{#1}
\providecommand{\decosix}{}
\providecommand{\starredbullet}{\ensuremath{\bullet}}
\pagestyle{empty}
\setlength{\parindent}{0pt}
\usepackage{adjustbox}\DeclareUnicodeCharacter{2236}{:}\DeclareUnicodeCharacter{2190}{\ensuremath{\leftarrow}}
\providecommand{\R}{\mathbb{R}}
\providecommand{\N}{\mathbb{N}}
\providecommand{\Z}{\mathbb{Z}}
\providecommand{\Q}{\mathbb{Q}}
\providecommand{\D}{\mathbb{D}}
\providecommand{\C}{\mathbb{C}}
\providecommand{\e}{{\mathrm{e}}}
\providecommand{\vect}[1]{\overrightarrow{#1}}
\providecommand{\vectt}[1]{\overrightarrow{\mathrm{#1}}}
\providecommand{\vvt}[1]{\overrightarrow{\mathrm{#1}}}
\providecommand{\np}[2][]{{#2~#1}}
\providecommand{\vv}[1]{{\overrightarrow{#1}}}
\renewcommand{\d}{{\mathrm{d}}}
\providecommand{\degres}{{^\circ}}
\providecommand{\degre}{{^\circ}}
\providecommand{\blue}{\color{blue}}
\providecommand{\red}{\color{red}}
\providecommand{\enskip}{\quad}
\providecommand{\llbracket}{\lbrack\!\lbrack}
\providecommand{\rrbracket}{\rbrack\!\rbrack}
\providecommand{\euro}{{\text{€}}}
\providecommand{\no}{{\text{n°}}}
\providecommand{\interval}[2]{{[#1;#2]}}
\providecommand{\intervFO}[2]{{[#1;#2[}}
\providecommand{\dotfill}{\ldots}

\providecommand{\ds}{\displaystyle}
\providecommand{\textding}[1]{#1}
\providecommand{\Oij}{\(\mathrm{O};\vec{\imath},\vec{\jmath}\)}
\providecommand{\Oijk}{\(\mathrm{O};\vec{\imath},\vec{\jmath},\vec{k}\)}
\providecommand{\Ouv}{\(\mathrm{O};\vec{u},\vec{v}\)}
\providecommand{\intervFF}[2]{\ensuremath{[#1;#2]}}
\renewcommand{\euro}{\ifmmode\text{\texteuro}\else\texteuro\fi}
\DeclareUnicodeCharacter{00B0}{\ifmmode^{\circ}\else\textdegree\fi}
\DeclareUnicodeCharacter{20AC}{\ifmmode\text{\texteuro}\else\texteuro\fi}
% Énoncés : Éducation nationale. Transcription et corrigés : APMEP, Denis Vergès et contributeurs. Découpage et classement : CoopMaths, Sébastien Lozano, Éric Elter et contributeurs. Adaptation de la présentation : Maths Explorer.
% CC BY-SA : https://coopmaths.fr/www/about/
% https://coopmaths.fr/alea/static/bac/2021/tex/bac_2021_06_metropole_8_cor.tex
\begin{document}
\noindent\textit{APMEP / CoopMaths. CC BY-SA. Présentation adaptée par Maths Explorer.}\par\medskip
\textbf{Commun à tous les candidats}

\medskip

Soit la suite $\left(u_n\right)$ définie par : $u_0 = 1$ et, pour $n$,
$ u_{n+1}  = \dfrac{4u_n}{u_n + 4}.$

\smallskip

\begin{enumerate}
\item


























On peut conjecturer que quel que soit $n$, \, $\dfrac{4}{u_n} = n + 4$.
\item On veut démontrer par récurrence que, pour tout entier naturel $n$, on a : $u_n > 0$.

\emph{Initialisation} : $u_0 = 1 > 0$ : la proposition est vraie au rang $0$.

\emph{Hérédité} : supposons qu'il existe $n \in \N$ tel que $u_n > 0$.

Alors $u_n + 4  >4 > 0$ donc : $ \dfrac{1}{u_n + 4} >0$ et donc $ \dfrac{4}{u_n + 4} >0$

Or $u_n>0$ (hypothèse de récurrence) donc $\dfrac{4u_n}{u_n + 4}>0$.

Soit finalement : $u_{n+1} >0$.; la proposition est vraie au rang $n + 1$.

La proposition est vraie au rang $0$ et si elle est vraie au rang $n \in \N$, elle est vraie au rang $n + 1$.

D'après le principe de récurrence on a démontré que pour tout $n \in \N$, \, $u_n > 0$.
\item 
Quel que soit $n \in \N$, \, $u_{n+1} - u_n = \dfrac{4u_n}{u_n + 4} - u_n = \dfrac{4u_n - u_n^2 - 4u_n}{u_n + 4} = - \dfrac{u_n^2}{u_n + 4}$.

Comme $u_n + 4 > 4 > 0$ d'où l'inverse $\dfrac{1}{u_n + 4} > 0$ et comme $u_n^2 > 0$, $\dfrac{u_n^2}{u_n + 4} > 0$ et finalement 

$- \dfrac{u_n^2}{u_n + 4} < 0$.

On a donc quel que soit $n \in \N$, \, $u_{n+1} - u_n < 0$ : la suite $\left(u_n\right)$ est décroissante.
\item 
La suite $\left(u_n\right)$ est décroissante et minorée par $0$ : d'après le théorème de la convergence monotone, la suite est convergente vers un réel $\ell \geqslant 0$

\item 






On a pour $n \in \N$, \\$v_{n+1} - v_n = \dfrac{4}{u_{n+1}} - \dfrac{4}{u_n} = \dfrac{4}{\frac{4u_n}{u_n + 4}} - \dfrac{4}{u_n} = \dfrac{4\left(u_n + 4 \right)}{4u_n} - \dfrac{4}{u_n} = \dfrac{u_n + 4}{u_n} - \dfrac{4}{u_n} = \dfrac{u_n + 4 - 4}{u_n} = \dfrac{u_n}{u_n} = 1$.

$v_{n+1} - v_n = 1$, quel que soit $n \in \N$ montre que la suite $\left(v_n\right)$ est une suite arithmétique de raison $r = 1$.

Son premier terme est $v_0 = \dfrac{4}{u_0} = \dfrac{4}{1} = 4$.

\item 
On sait qu'alors pour $n \in \N$, \, $v_n = v_0 + nr = 4 + n$.

Or $v_n = \dfrac{4}{u_n} \iff u_n = \dfrac{4}{v_n}$ donc $u_n = \dfrac{4}{4 + n}$ quel que soit $n \in \N$.

$u_n = \dfrac{1}{1 + \frac{n}{4}}$, donc comme $\displaystyle\lim_{n \to + \infty} 1 + \frac{n}{4} = + \infty$, \, $\displaystyle\lim_{n \to + \infty} u_n = 0$.


La limite de la suite $\left(u_n\right)$ est donc $0$.
\end{enumerate}

\bigskip

\end{document}
